%====================
% EXPERIENCE A
%====================
\subsection{{Data Science Director \hfill July 2025--- Present}}
\subtext{AstraZeneca \hfill Gaithersburg, MD}
\begin{zitemize}
    \item Lead statistical thinking and decision support for Oncology R\&D precision diagnostics across lung, bladder, colorectal, gastric, and hematologic cancer programs
    \item Quantify benefit, risk, and uncertainty of ctDNA diagnostic strategies across 20+ clinical programs at various development stages, providing design options that enable informed go/no-go decisions
    \item Direct statistical staff and CRO partners on study-level and program-level deliverables, ensuring quality and on-time delivery
    \item Developed Bayesian hierarchical models and simulation frameworks for ctDNA assay harmonization across 15+ studies, enabling cross-platform evidence synthesis for regulatory submission support
    \item Lead communities of practice in Bayesian methods and R (1,500+ members); mentor and coach statistical staff through workshops and training programs
\end{zitemize}

%====================
% EXPERIENCE B
%====================
\subsection{{Data Science Associate Director \hfill Aug 2020 --- June 2025}}
\subtext{AstraZeneca \hfill Gaithersburg, MD}
\begin{zitemize}
    \item Led statistical design and analysis for companion diagnostic development from exploration through Phase III trials, serving as lead statistician across multiple indications
    \item Designed Bayesian frameworks for diagnostic sample sizing and target product profiles, directly informing Phase III trial strategy and regulatory interactions
    \item Led multi-vendor diagnostic bakeoffs, developing novel statistical approaches to evaluate assay performance and inform assay selection for pivotal trials
    \item Delivered production ML pipelines for cell type annotation from scRNA-seq data; co-authored 9 peer-reviewed papers communicating statistical findings to broader scientific community
    \item Directed CRO partnerships for statistical programming and data deliverables across multiple concurrent studies
\end{zitemize}

%====================
% EXPERIENCE C
%====================
\subsection{{Consultant Statistician \hfill Jan 2010 --- Aug 2020}}
\subtext{National Institute of Arthritis, Musculoskeletal and Skin Disorders, NIH\hfill Bethesda, MD}
\begin{zitemize}
    \item Lead statistician for rheumatology clinical and observational studies, independently performing and directing all statistical work across 20+ projects
    \item Provided high-quality decision support through study design, statistical modeling, Bayesian meta-analyses, causal inference, and simulation
    \item Investigated and applied novel statistical approaches including spatial epidemiology, adaptive methods, and complex survival models
    \item Delivered 40+ peer-reviewed publications; built strong cross-functional relationships with clinical investigators, epidemiologists, and program leaders
\end{zitemize}

%====================
% EXPERIENCE D
%====================
\subsection{{Assistant Professor \hfill Mar 2006 --- Aug 2009}}
\subtext{Thomas Jefferson University \hfill Philadelphia, PA}
\begin{zitemize}
    \item Led statistical collaboration for cancer and cardiology research programs, producing 25+ peer-reviewed publications
    \item Developed novel statistical approaches for high-dimensional genomic data (microarray analyses) across multiple cancer studies
    \item Mentored graduate students and junior researchers in biostatistical methods
\end{zitemize}

\subsection{Co-authored 88 peer-reviewed papers (\href{https://scholar.google.com/citations?hl=en&user=EgaGUCwAAAAJ}{Google Scholar})}
